\clearpage
\begin{trapbox}
\textbf{误区一：横向抛物线的对称弦竖直，仍直接写 $1/k_{AB}$。}
\begin{center}% p=1, free coordinates +1,-1, reflected with identical scale.
\begin{tikzpicture}[font=\small,>=stealth]
\draw[gray!45,rounded corners=3pt] (0,0) rectangle (14.2,10.0);
\node at (7.1,9.45) {对称的两个端点：竖直弦与水平弦的斜率含义不同};
\node at (3.55,8.45) {$y^2=2x$};
\node at (10.65,8.45) {$x^2=2y$};
\begin{scope}[shift={(2.6,5.1)},scale=1.25]
\draw[->,gray] (-1.5,0)--(2.1,0) node[above] {$x$};
\draw[->,gray] (0,-1.5)--(0,2.1) node[left] {$y$};
\draw[teal!65!black,thick,samples=90,domain=-1.4:1.8,variable=\t] plot ({\t*\t/2},\t);
\draw[blue!75!black,thick] (0,0)--(.5,1);
\draw[orange!85!black,thick] (0,0)--(.5,-1);
\draw[green!45!black,very thick] (.5,1)--(.5,-1);
\fill (0,0) circle(2pt) node[below left=7pt] {$O$};
\fill (.5,1) circle(2pt) node[right=9pt] {$A$};
\fill (.5,-1) circle(2pt) node[right=9pt] {$B$};
\end{scope}
\begin{scope}[shift={(9.7,5.1)},scale=1.25]
\draw[->,gray] (-1.5,0)--(2.1,0) node[above] {$x$};
\draw[->,gray] (0,-1.5)--(0,2.1) node[left] {$y$};
\draw[teal!65!black,thick,samples=90,domain=-1.4:1.8,variable=\t] plot (\t,{\t*\t/2});
\draw[blue!75!black,thick] (0,0)--(1,.5);
\draw[orange!85!black,thick] (0,0)--(-1,.5);
\draw[green!45!black,very thick] (-1,.5)--(1,.5);
\fill (0,0) circle(2pt) node[below left=7pt] {$O$};
\fill (1,.5) circle(2pt) node[above=14pt] {$A$};
\fill (-1,.5) circle(2pt) node[above=14pt] {$B$};
\end{scope}
\node at (3.55,2.6) {$AB:\ x=\frac12$，$k_{AB}$ 不存在};
\node at (10.65,2.6) {$AB:\ y=\frac12$，$k_{AB}=0$};
\node at (3.55,1.55) {$\displaystyle\frac{x_B-x_A}{y_B-y_A}=0$};
\node at (10.65,1.55) {$\displaystyle k_{OA}+k_{OB}=\frac12-\frac12=0$};
\node at (7.1,.55) {同一尺度下，横向模型的竖直弦关于 $y=x$ 翻转后成为纵向模型的水平弦。};
\end{tikzpicture}
\end{center}

\textbf{错法：}由 $k_{OA}=2$、$k_{OB}=-2$，算出 $1/k_{AB}=0$。

\textbf{为什么错：}竖直直线的斜率不存在，因而它的倒数也不能按普通实数定义。\textbf{正确写法}是
\[
\frac1{k_{OA}}+\frac1{k_{OB}}
=\frac{x_B-x_A}{y_B-y_A}=0,
\]
再说明 $AB$ 竖直。不能把“不存在的斜率”当作“无穷大的斜率”。

\medskip
\textbf{误区二：把两种开口方向的加法直接混用。}

\textbf{错法：}把 $y^2=2px$ 写成 $k_{OA}+k_{OB}=k_{AB}$。

\textbf{正确区分：}横向开口加\textbf{倒斜率}，纵向开口加\textbf{斜率}。因为交换横纵坐标会让通常的非零有限斜率变成倒数。
\end{trapbox}

\clearpage
\begin{trapbox}
\textbf{误区三：误以为所有斜率都必须非零。}

纵向开口 $x^2=2py$ 的 $AB$ 可以水平，此时 $k_{AB}=0$ 合法，并不妨碍 $k_{OA}+k_{OB}=k_{AB}$。横向开口的 $OA,OB$ 在 $A,B\ne O$ 的前提下，斜率必非零。

\medskip
\textbf{误区四：允许 $A=O$ 或 $B=O$，却还讨论 $OA$ 或 $OB$。}
\begin{center}% Identical coordinate ranges; y^2=2x in both panels.
\begin{tikzpicture}[font=\small,>=stealth]
\draw[gray!45,rounded corners=3pt] (0,0) rectangle (14.2,9.5);
\node at (7.1,8.95) {先检查点是否不同、是否离开顶点，再讨论三条边};
\node at (3.55,7.95) {$A=O$};
\node at (10.65,7.95) {$A=B$};
\foreach \dx in {1.8,8.9}{
 \begin{scope}[shift={(\dx,3.1)},scale=1.15]
 \draw[->,gray] (-.6,0)--(3.1,0) node[above] {$x$};
 \draw[->,gray] (0,-.6)--(0,3.1) node[left] {$y$};
 \draw[teal!65!black,thick,samples=100,domain=-.7:2.35,variable=\t] plot ({\t*\t/2},\t);
 \draw[blue!75!black,very thick] (0,0)--(2,2);
 \end{scope}
}
\fill (1.8,3.1) circle(2.5pt) node[below right=8pt] {$A=O$};
\fill (4.1,5.4) circle(2.5pt) node[above left=8pt] {$B$};
\fill (8.9,3.1) circle(2.5pt) node[below right=8pt] {$O$};
\fill (11.2,5.4) circle(2.5pt) node[above left=8pt] {$A=B$};
\node[align=center] at (3.55,1.25) {两个重合点 $O,A$\\[8pt]不能确定直线 $OA$};
\node[align=center] at (10.65,1.25) {两个端点重合\\[8pt]不能确定弦 $AB$};
\end{tikzpicture}
\end{center}

若 $A=O$，那么“直线 $OA$”由两个重合点给出，方向无法确定；所以本结论必须写明 $A,B$ 均不是顶点 $O$，且 $A\ne B$。

\medskip
\textbf{误区五：认为 $p$ 必须大于 $0$。}

对标准形式 $y^2=2px$、$x^2=2py$，$p<0$ 分别对应向左、向下开口，公式不变。注意 $p\ne0$ 是必要条件；不要把退化方程当成抛物线。
\end{trapbox}

\clearpage
\begin{trapbox}
\textbf{误区六：平方差约分时没有检查两点是否不同。}

\textbf{错法：}从
\[
\frac{b^2-a^2}{b-a}=a+b
\]
直接约掉 $b-a$，却没有说明 $a\ne b$。

\textbf{怎样改：}抛物线按本节参数表示时，同一个参数对应唯一一点；既然 $A\ne B$，必有 $a\ne b$，这才可以约分。

\medskip
\textbf{误区七：只背结论，不知道从哪里找斜率之和。}

遇到横向开口，写 $x=y^2/(2p)$，把两点的\textbf{纵坐标}设成 $a,b$；遇到纵向开口，写 $y=x^2/(2p)$，把两点的\textbf{横坐标}设成 $a,b$。两种情形都能一眼得到
\[
\boxed{\frac{a+b}{2p}}.
\]
这是整个二级结论最值得迁移到其他题目的构造方式。
\end{trapbox}

\clearpage
\begin{trapbox}
\textbf{误区八：把顶点连线、端点切线和弦当成同一种直线。}
\begin{center}% x^2=2y, p=1; A=(-1,1/2), B=(2,2); slopes -1,2,1/2.
\begin{tikzpicture}[font=\small,>=stealth]
\draw[gray!45,rounded corners=3pt] (0,0) rectangle (14.2,10.0);
\node at (7.1,9.45) {纵向：两端切线的斜率取平均，得到弦的斜率};
\begin{scope}[shift={(3.8,3.1)},scale=1.15]
\draw[->,gray] (-2.4,0)--(3.4,0) node[above] {$x$};
\draw[->,gray] (0,-1.3)--(0,4.3) node[left] {$y$};
\draw[teal!65!black,thick,samples=100,domain=-2.2:2.8,variable=\t] plot (\t,{\t*\t/2});
\draw[blue!75!black,dashed,thick] (0,0)--(-1,.5);
\draw[orange!85!black,dashed,thick] (0,0)--(2,2);
\draw[blue!75!black,very thick] (-2.2,1.7)--(.6,-1.1);
\draw[orange!85!black,very thick] (.6,-.8)--(2.8,3.6);
\draw[green!45!black,very thick] (-1,.5)--(2,2);
\fill (0,0) circle(2pt) node[below left=10pt] {$O$};
\fill (-1,.5) circle(2pt) node[above=14pt] {$A$};
\fill (2,2) circle(2pt) node[right=14pt] {$B$};
\node[blue!75!black] at (-1.6,2.1) {$t_A$};
\node[orange!85!black] at (3.25,3.25) {$t_B$};
\end{scope}
\node[anchor=west,align=left] at (8.9,5.7)
 {$x^2=2y,\quad p=1$\\[13pt]
 $A(-1,\frac12),\ B(2,2)$\\[15pt]
 $k_A=-1,\quad k_B=2$\\[15pt]
 $k_{AB}=\dfrac{-1+2}{2}=\dfrac12$};
\node at (7.1,1.3) {虚线是顶点连线，实线 $t_A,t_B$ 是切线；二者不能混淆。};
\node at (7.1,.6) {$k_{OA}=-\frac12,\ k_{OB}=1$，切线斜率分别是顶点连线斜率的 $2$ 倍。};
\end{tikzpicture}
\end{center}
\textbf{错法：}图中切线斜率为 $-1,2$，直接相加得到 $k_{AB}=1$。
\textbf{为什么错：}相加的是虚线 $OA,OB$ 的斜率，切线斜率应先取平均。
\textbf{正确替代：}$k_{AB}=(-1+2)/2=1/2$；也可用顶点连线斜率 $-1/2,1$ 相加。

\medskip
\textbf{误区九：切线公式不分开口方向，或忘记斜率的倒数。}
\textbf{错法：}横向切线斜率为 $-1,1/2$ 时，写 $k_{AB}=(-1+1/2)/2$。
\textbf{正确替代：}横向取倒斜率的平均，$1/k_{AB}=(-1+2)/2=1/2$，故 $k_{AB}=2$。纵向才是普通斜率的平均。

\medskip
\textbf{误区十：倒数和为零，仍强行取倒数。}
\textbf{错法：}把 $uv/(u+v)$ 的分母 $0$ 当作一个能代入的数。
\textbf{正确替代：}横向顶点连线斜率和为 $0$ 时，改用坐标差或加强式判断竖直弦；横向切线倒斜率平均值为 $0$ 时也先判断竖直弦，不写 $1/k_{AB}=0$。
\end{trapbox}

\clearpage
\begin{trapbox}
\textbf{误区十一：普通弦也套焦点弦长度倒数和。}
\begin{center}% p=2; focal chord x=1 vs nonfocal chord x=4, identical scale.
\begin{tikzpicture}[font=\small,>=stealth]
\draw[gray!45,rounded corners=3pt] (0,0) rectangle (14.2,10.5);
\node at (7.1,9.95) {只改变弦的位置：经过焦点这个条件不能省};
\node at (3.55,9.05) {$AB:x=1$，经过 $F$};
\node at (10.65,9.05) {$AB:x=4$，不经过 $F$};
\foreach \offset in {1.5,8.6}{
\begin{scope}[shift={(\offset,5.0)},scale=.7]
\draw[->,gray] (-.5,0)--(5.7,0) node[above] {$x$};
\draw[->,gray] (0,-4.6)--(0,4.6) node[left] {$y$};
\draw[teal!65!black,thick,samples=100,domain=-4.4:4.4,variable=\t] plot ({\t*\t/4},\t);
\end{scope}}
\begin{scope}[shift={(1.5,5.0)},scale=.7]
\draw[green!45!black,very thick] (1,-2)--(1,2);
\fill (1,2) circle(2pt) node[right=12pt] {$A$};
\fill (1,-2) circle(2pt) node[right=12pt] {$B$};
\fill (1,0) circle(2pt) node[below right=12pt] {$F$};
\end{scope}
\begin{scope}[shift={(8.6,5.0)},scale=.7]
\draw[blue!75!black,dashed,thick] (1,0)--(4,4);
\draw[orange!85!black,dashed,thick] (1,0)--(4,-4);
\draw[green!45!black,very thick] (4,-4)--(4,4);
\fill (4,4) circle(2pt) node[right=12pt] {$A$};
\fill (4,-4) circle(2pt) node[right=12pt] {$B$};
\fill (1,0) circle(2pt) node[above=14pt] {$F$};
\end{scope}
\node at (3.55,1.55) {$FA=FB=2,\quad \frac1{FA}+\frac1{FB}=1$};
\node at (10.65,1.55) {$FA=FB=5,\quad \frac1{FA}+\frac1{FB}=\frac25$};
\node at (7.1,.55) {两图同尺度：$y^2=4x,\ F(1,0)$；普通弦不能套用焦点弦定值。};
\end{tikzpicture}
\end{center}
\textbf{错法：}只看到两个端点在 $y^2=4x$ 上，就写 $1/FA+1/FB=1$。
\textbf{为什么错：}右图 $AB:x=4$ 不经过 $F(1,0)$，两段长均为 $5$。
\textbf{正确替代：}先验证弦过焦点且有两个不同有限端点；否则直接求距离，右图倒数和是 $2/5$。

\medskip
\textbf{误区十二：长度倒数、斜率倒数、长度平方倒数混用。}
\textbf{错法：}把 $FA$ 换成 $k_{OA}$，或把 $1/h^2$ 改成 $1/h$。
\textbf{为什么错：}这些符号表示不同对象，推导也不同。
\textbf{正确替代：}先写清对象，再选择平方差、焦点距离关系，或面积加勾股；相似外形不能替代条件检查。

\medskip
\textbf{误区十三：平移后的抛物线仍连接原点。}
\textbf{错法：}对 $y=(x-2)^2+1$ 仍计算 $k_{OA}+k_{OB}$。
\textbf{正确替代：}连接实际顶点 $V(2,1)$；平移保持的是 $VA,VB,AB$ 的方向，不是换成任意定点后的方向。
\end{trapbox}

\clearpage
\begin{trapbox}
\textbf{误区十四：直线只有一个交点，就一定是切线。}
\textbf{错法：}横向抛物线 $y^2=2px$ 与 $y=1$ 只有一个交点，所以 $y=1$ 是切线。
\textbf{为什么错：}代入 $y=1$ 后方程降为关于 $x$ 的一次方程；这与二次方程出现重根不是一回事。
\textbf{正确替代：}使用 $x=qy+d$ 联立，检查关于 $y$ 的二次方程是否有重根，再比较系数求切线。顶点处切线为 $x=0$，普通斜率不存在。

\medskip
\textbf{最后检查：}本节 $p$ 在主斜率公式中允许正负；抛物线焦点弦长度式 $2/p$ 则按 $y^2=2px,\ p>0$ 的条件使用。
椭圆式的 $a,b$ 是半轴长，直角三角形式的 $a,b$ 是直角边长，不可在不同场景间直接替换。
\end{trapbox}
